\chapter{Introduction}
\section{Background and motivation}
Transformer language models generate text one token at a time. Before producing each new token, every attention layer compares the current position with all earlier positions of the conversation~\cite{vaswani2017attention}. To avoid recomputing the representations of earlier tokens at every step, the model stores, for every layer and attention head, a \emph{key} vector and a \emph{value} vector for each token it has processed. This growing store is the key--value (KV) cache. It is created during a conversation, it changes with every token that is read or written, and it is discarded when the conversation ends, while the trained weights of the model never change. In that sense the KV cache is the short-term, runtime memory of the model: everything the model ``remembers'' about the current conversation is held there.

Research on the KV cache has so far concentrated on its cost. Long conversations make the cache large, so a substantial body of work studies how to evict, compress, stream or share it efficiently~\cite{xiao2024sinks,zhang2023h2o,liu2023scissorhands,oren2024msrnn,kwon2023vllm}. Several of these studies observe that the cache is highly structured: a few early ``attention sink'' tokens receive large attention regardless of their meaning~\cite{xiao2024sinks}, and a small set of ``heavy-hitter'' entries dominates attention throughout generation~\cite{zhang2023h2o}. However, these structures are studied only to decide which entries can be safely discarded. Whether the information held in the cache shapes what the model \emph{believes} during a conversation, and in particular whether it carries and amplifies social bias, has not been examined.

The question has become practical because language models now work in teams. Multi-agent frameworks let one model search, another write and a third check the result~\cite{wu2023autogen,du2024debate}, and serving systems already reuse cached keys and values across requests to save computation~\cite{kwon2023vllm,yao2025cacheblend}. When agents exchange text, a human can in principle read what was said. When they exchange cache entries, the transferred information is a set of numerical vectors that no person can read. If a one-sided idea can travel from one agent to another through such shared memory, it could spread through an AI system without leaving any visible trace in the exchanged text. A recent survey of bias and fairness in language models~\cite{gallegos2024bias} shows that most mitigation methods act before or during training, which is expensive and out of reach for small institutions, and that few methods act on the model at inference time.

\textbf{Relevance to the Sustainable Development Goals.} The project was proposed under United Nations Sustainable Development Goal~10, \emph{Reduced Inequalities}~\cite{un2015agenda}. Language models are increasingly used to shortlist job applicants, answer health questions, give educational feedback and run public-service chatbots. When these systems carry social bias, the harm falls mostly on groups that are already disadvantaged by gender, religion, disability, origin or economic status. This directly concerns Target~10.2, which calls for the social, economic and political inclusion of all people irrespective of such characteristics, and Target~10.3, which calls for equal opportunity and the elimination of discriminatory practices. The project contributes by making it possible to observe where bias is held inside a running model and to test whether it can be removed without retraining, which lowers the cost of auditing for the smaller organisations that most need it. The work also supports two secondary goals. Under SDG~9, \emph{Industry, Innovation and Infrastructure}, and its Target~9.5 on enhancing scientific research and technological capability, it delivers open, reproducible research software that runs on commodity hardware. Under SDG~16, \emph{Peace, Justice and Strong Institutions}, and its Target~16.6 on effective, accountable and transparent institutions, it provides tools for inspecting and documenting the behaviour of automated decision systems. The intended beneficiaries are the people affected by automated decisions, and the researchers, auditors and developers who build or regulate such systems. The contribution demonstrated in this report is a measurement and intervention toolkit evaluated in a controlled setting; its benefit for real deployments is an expected outcome that depends on the larger-scale validation described in Chapter~7.

\section{Objectives}
The project has five objectives, taken from the approved design document and stated here together with the evaluation method and deliverable for each.

The first objective (O1) is to build a runtime-memory instrumentation framework that records every KV cache entry of a frozen language model together with its provenance: layer, attention head, absolute position, conversation turn, agent and source channel. It is evaluated by unit tests that check provenance stamping, cache transplantation and deterministic replay, and its deliverable is an installable Python package with a trace store.

The second objective (O2) is to detect and characterise persistent \emph{memory attractors}, that is, cache entries that keep receiving attention long after they were written, and to separate them from the known attention-sink phenomenon. It is evaluated with four metrics (reuse, lifetime, persistence and centrality) under pre-registered thresholds, and its deliverable is the attractor statistics for every conversation.

The third objective (O3) is to test whether bias is concentrated in identifiable \emph{bias reservoirs}, specific attention heads whose cached values predict the bias of a conversation. It is evaluated by cross-validated probe AUC with false-discovery-rate control and by the measured effect of ablating each head, and its deliverable is a ranked reservoir map.

The fourth objective (O4) is to measure \emph{memory contagion}, the part of bias transfer between two agents that is carried by shared cache entries rather than by exchanged text, and to test whether bias rises abruptly at a critical level of one-sided input (a phase transition). It is evaluated by a three-arm experiment with bootstrap confidence intervals and a Mann--Whitney test, and by comparing logistic and linear fits, and its deliverable is the contagion decomposition and phase analysis.

The fifth objective (O5) is to change behaviour by editing the cache alone. Ablation, perturbation and gating interventions are tested under seed-matched replay against random-head and sink-only negative controls, and an intervention is accepted only if it stays within a coherence budget of at most 5\% perplexity increase and at most 2 points of factual-accuracy loss. Its deliverable is a causal intervention toolkit with the accept or reject decision for each intervention.

\section{Scope and constraints}
The intended users of the framework are researchers, auditors and developers who need to inspect how a language model uses its runtime memory. The project covers decoder-only transformer models whose weights are publicly available, two-agent conversations between instances of the same model, and social bias measured on fictional groups, with factual question answering used as a coherence check. The model weights are never modified: every measurement and every intervention acts only on the KV cache during inference.

Three constraints shaped the work. First, the available hardware was a laptop with a 4~GB NVIDIA GTX 1650~Ti GPU. The design document targets 7--8 billion parameter models such as Llama-3-8B, which need roughly 16~GB of GPU memory for their weights alone. The framework supports these models through configuration files, but in this phase it was evaluated on a controlled reference transformer and on the 0.5 billion parameter Qwen2.5-0.5B-Instruct model~\cite{qwen2024}. Second, to make the direction of bias known by construction and to avoid making statements about real communities, the probes use invented group names such as ``Vorlan'' and ``Tessic'' rather than real demographic groups, following the question formats of established benchmarks~\cite{parrish2022bbq,nadeem2021stereoset,nangia2020crows,lin2022truthfulqa}. Third, cache sharing between agents is only defined for agents running the same base model, because cached keys and values are not portable across tokenizers or architectures. Fine-tuning, retraining and debiasing of weights are outside the scope of the project.

\section{Project report organization}
The rest of this report is organised as follows. Chapter~2 states the problem and analyses the limitations of existing approaches. Chapter~3 reviews twelve research papers on KV cache structure, interpretability and intervention, multi-agent interaction and bias measurement, and compares them in a table. Chapter~4 lists the functional, non-functional, hardware and software requirements. Chapter~5 presents the architecture, the experimental protocol, the metrics and the intervention procedure, together with the demonstration website. Chapter~6 describes the evaluation setup and criteria, reports the results for both models, compares them with existing work and with the expected outcomes of the design document, and discusses limitations. Chapter~7 concludes and lays out the work planned for the next phase. Appendix~A documents how to install, run and verify the software.
